%% ma: L12-B16-0020
%% lop: 12
%% chuong: 5
%% bai: 16
%% dang: 12.16.13
%% loai: TN4
%% muc: VD
%% dap_an: "D"
%% nguon: "Ảnh giáo viên gửi 04/11/2026 (Huỳnh Văn Ánh, Chuyên đề V - Phương pháp tọa độ trong không gian), Đề test theo bài, Phần I Câu 11"
%% kien_thuc: "Bài 16 (góc giữa hai đường thẳng), Bài 7 và Bài 8 (tọa độ hóa hình chóp đều, tích vô hướng), Lớp 11 (hình chóp đều, đường cao của chóp đều)"
%% trang_thai: chinh-thuc
%% ghi_chu: "Đáp án gốc D khớp (sympy): cos = 219/căn 471161 = 0,3190..., góc khoảng 71,39 độ. Ảnh chụp kim tự tháp trong đề gốc là ảnh minh họa đời sống, không chứa dữ kiện toán nên bỏ; thay bằng hình chóp S.ABCD vẽ bằng TikZ (G nằm trên mặt (SCD) bị khuất nên BG vẽ nét đứt). Phân vân giữa VD và VDC, chọn mức thấp hơn. Đã duyệt 04/11/2026."
\begin{ex}
Kim tự tháp Kheops ở Ai Cập có dạng hình chóp $S.ABCD$, có đáy là hình vuông với cạnh dài $230$ m, các cạnh bên bằng nhau và dài $219$ m (theo britannica.com) (hình vẽ). Gọi $G$ là trọng tâm tam giác $SCD$. Tính góc giữa hai đường thẳng $BG$ và $SA$ (làm tròn tới hàng đơn vị độ).
\begin{center}
\begin{tikzpicture}[x={(1cm,0cm)},y={(0.8cm,0.5cm)},z={(0cm,1cm)},scale=0.9]
  \coordinate (A) at (0,0,0);
  \coordinate (B) at (4,0,0);
  \coordinate (C) at (4,4,0);
  \coordinate (D) at (0,4,0);
  \coordinate (S) at (2,2,2.55);
  \coordinate (G) at (2,3.3333,0.85);
  \draw[khuat] (D)--(A) (D)--(C) (D)--(S) (B)--(G);
  \draw (A)--(B)--(C) (S)--(A) (S)--(B) (S)--(C);
  \fill (G) circle (1.3pt);
  \node[above,font=\footnotesize] at (S) {$S$};
  \node[below left,font=\footnotesize] at (A) {$A$};
  \node[below right,font=\footnotesize] at (B) {$B$};
  \node[right,font=\footnotesize] at (C) {$C$};
  \node[above left,font=\footnotesize] at (D) {$D$};
  \node[right,font=\footnotesize] at (G) {$G$};
\end{tikzpicture}
\end{center}
\choice{$66^\circ$}{$69^\circ$}{$70^\circ$}{\True $71^\circ$}
\loigiai{Gọi $O=AC\cap BD$. Vì $S.ABCD$ là hình chóp đều nên $SO\perp(ABCD)$. Đặt $AB=230=a$, $SA=219=b$. Khi đó $OA=OB=OC=OD=\dfrac{a\sqrt2}2$ và $SO=\sqrt{SA^2-OA^2}=\dfrac{\sqrt{4b^2-2a^2}}2$.

Chọn hệ trục $Oxyz$ với $Ox$ qua $C$, $Oy$ qua $D$, $Oz$ qua $S$: $A\left(-\dfrac{a\sqrt2}2;0;0\right)$, $B\left(0;-\dfrac{a\sqrt2}2;0\right)$, $C\left(\dfrac{a\sqrt2}2;0;0\right)$, $D\left(0;\dfrac{a\sqrt2}2;0\right)$, $S\left(0;0;\dfrac{\sqrt{4b^2-2a^2}}2\right)$.

Vì $G$ là trọng tâm tam giác $SCD$ nên $G\left(\dfrac{a\sqrt2}6;\dfrac{a\sqrt2}6;\dfrac{\sqrt{4b^2-2a^2}}6\right)$.

Ta có $\overrightarrow{SA}=\left(-\dfrac{a\sqrt2}2;0;-\dfrac{\sqrt{4b^2-2a^2}}2\right)$, $\overrightarrow{BG}=\left(\dfrac{a\sqrt2}6;\dfrac{2a\sqrt2}3;\dfrac{\sqrt{4b^2-2a^2}}6\right)$.

$|\overrightarrow{SA}\cdot\overrightarrow{BG}|=\dfrac{b^2}3$, $|\overrightarrow{SA}|=b$, $|\overrightarrow{BG}|=\dfrac{\sqrt{b^2+8a^2}}3$. Suy ra $\cos(SA,BG)=\dfrac{b}{\sqrt{b^2+8a^2}}=\dfrac{219}{\sqrt{471161}}\approx0{,}3191$.

Vậy $(SA,BG)\approx71^\circ$. Chọn D.}
\end{ex}
